% \documentclass[addpoints]{exam}
\documentclass[addpoints,12pt]{exam}

\usepackage{arev}

\usepackage[slovene]{babel}
\usepackage{amsfonts}
\usepackage{amssymb}
\usepackage{float}
\usepackage{graphicx}
\usepackage{enumitem}
\usepackage{color}
\usepackage{eurosym}

%Komentar naslednje vrstice glede na verzijo testa -- rešitve ali prostor za odgovore:
% \printanswers

% \linespread{2}


\pointsinrightmargin
\marginbonuspointname{}


\renewcommand{\solutiontitle}{\noindent\textbf{Rešitve in točkovanje:}\par\noindent}

\colorgrids
\definecolor{GridColor}{gray}{0.7}
\setlength{\gridsize}{7mm}

\extrawidth{5mm}
\setlength{\rightpointsmargin}{1.5cm}

\begin{document}

\runningfooter{}{}{}

\begin{center}
    \fbox{\fbox{\parbox{15cm}{\centering
    Prvo pisno ocenjevanja znanja \\ 1. letnik -- test A \\ 23. oktober 2024 \\ (Osnove logike in teorije množic, Naravna števila)}}}
\end{center}

\vspace{0.1in}
\makebox[\textwidth]{Ime in priimek, oddelek:\enspace\hrulefill}


\begin{center}
    \chqword{Naloga}
    \chpword{Možne točke}
    \chbpword{Dodatne točke}
    \chsword{Dosežene točke}
    \chtword{Skupaj}  
    \combinedgradetable[h][questions]
\end{center}

\vspace{0.1in}


\begin{questions}

    % 1. vprašanje
    
    \question[4]
    Določite pravilnost izjave $(\lnot A\Rightarrow  B)\land(A\Leftrightarrow \lnot B)$ glede na izbor izjav $A$ in $B$.

    \begin{solution}[8cm]
???
    \end{solution}



\newpage
    % 2. vprašanje
    \question[3]
    Določite resničnostno vrednost izjav:
    \begin{parts}
        \part 
        Število $23$ je večkratnik števila 13 ali število $9$ je sodo praštevilo.\\
        \part 
        Če je število $2$ praštevilo, potem je praštevil neskončno mnogo.\\
        \part
        Ni res, da je ostanek pri deljenju s $4$ lahko samo $1$, $2$ ali $3$, in število $8$ je večje od števila $2$.\\
    \end{parts}

    \begin{solution}[1cm]
???
    \end{solution}




    % 3. vprašanje
    \question
    Dani sta množici $\mathcal{A}=\left\{a, b, c, h\right\}$ in $\mathcal{B}=\left\{b, c, h, i\right\}$.
    \begin{parts}
        \part[2]
        Ali velja $\mathcal{A}\nsubseteq\mathcal{B}$? Poiščite največjo možno množico $\mathcal{C}$, da bosta množici $\mathcal{A}$ in $\mathcal{B}$ njeni nadmnožici.
        \part[3]
        Zapišite $\mathcal{PB}$. Izračunajte koliko elementov vsebuje.
        \part[2]
        Zapišite in grafično predstavite kartezični produkt $\mathcal{B}\times\mathcal{A}$.
    \end{parts}

    % \begin{description}
    %     \item[a] $(x-a)^2+(y-b)^2-c^2=0$ 
    %     \item[b] $d^2(x-a)^2+f^2(y-b)^2-(df)^2=0$
    %     \item[c] $S(a,b)$
    %     \item[d] $e^2=d^2-f^2;\ d>f$
    % \end{description}
    % Krožnica: \fillin[a c][2cm]
    % \\ Elipsa:  \fillin[b c d][2cm]

    \begin{solution}[9cm]
???
    \end{solution}


    \newpage

    % 4. vprašanje
    \question
    Množica $\mathcal{A}$ je prava podmnožica množice $\mathcal{B}$, prav tako je $\mathcal{B}$ prava podmnožica množice $\mathcal{C}$.
    \begin{parts}
        \part[1]
        Grafično predstavite množico $(\mathcal{C}\setminus\mathcal{B})\cup\mathcal{A}$.
        \part[3]
        Izračunajte $\mathcal{B}\setminus\mathcal{C}$, $\mathcal{A}\cap\mathcal{C}$ in $\mathcal{A}\cup\mathcal{B}$.
    \end{parts}

    \begin{solution}[1cm]
???
    \end{solution}




    \newpage

    % 5. vprašanje
    \question
    Naj bo $\mathcal{U}=\mathbb{N}_{23}$. Dane so množice $\mathcal{A}=\left\{x;\ x=2n \land 3\leq n<10\right\}$, \\ $\mathcal{B}=\left\{x;\ x=3k-1 \land 3\leq k\leq 7\right\}$ in $\mathcal{C}=\left\{x;\ x=5m \land 1<m\leq 4\right\}$. 
    \begin{parts}
        \part[5]
        Določite naslednje množice: $\mathcal{A}\cap\mathcal{C}$, $\mathcal{C}\setminus\mathcal{A}$, $(\mathcal{A}\cup\mathcal{C})\setminus\mathcal{B}$ in $(\mathcal{A}\cup\mathcal{B})^\complement$.
        \part[1]
        Koliko elementov ima množica $\mathcal{A}\times\mathcal{C}$?
    \end{parts}

    \begin{solution}[9cm]
        ???
            \end{solution}


    \newpage
    % 6. vprašanje
    \question
    Pri uri športa so se dijaki in učitelj pogovarjali o smučeh.  Učitelj je dijake vprašal, katere vrste smuči imajo doma.
    $20$ dijakov ima \textit{alpske smuči}, $6$ jih ima \textit{turne smuči} in $17$ \textit{tekaške smuči}. Kar $12$ jih ima \textit{alpske} in \textit{tekaške smuči}, $4$ imajo \textit{turne} in \textit{alpske smuči},
    $3$ pa \textit{turne} in \textit{tekaške smuči}. En sam dijaki ima vse tri vrste smuči, en sam pa nima doma nobenih smuči. 
    \begin{parts}
        \part[1]
        Narišite diagram.
        \part[2]
        Koliko dijakov je bilo pri uri športa?
        \part[2]
        Koliko dijakov ima doma natanko dvoje vrst smuči?
    \end{parts}

    \begin{solution}[8cm]
        ???
            \end{solution}




    \newpage
        

    % 7. vprašanje
    \question[3]
    Na šoli je pet oddelkov s po $26$ dijaki, sedem oddelkov s po $27$ dijaki in štirje oddelki s po $28$ dijaki. 
    Med vsemi dijaki je $320$ deklet. Koliko fantov je na šoli? (Zapišite samo en račun.)

    \begin{solution}[11cm]
        ???
            \end{solution}



        % dodatno vprašanje
    
    \bonusquestion[3]
    \textit{Dodatna naloga:} Izjavo \textit{Obstaja takšno celo število $x$, da je deljivo s številom $3$ in s številom $2$.} zapišite s simboli in jo zanikajte v simbolnem zapisu.


\end{questions}



\end{document}